\section*{الفصل \ref{ch:b2:equilibrium-shifts} --- ثوابت التوازن وانزياحاته}

\begin{solution}{exo:b2:equilibrium-shifts:1}
$K^\circ = \exp(32\,800/(8.314 \times 298.15)) = \num{5.6e5}$، ومن أجل
$+\qty{32.8}{kJ/mol}$ مقلوبه، \num{1.8e-6}. ويتطلب العامل 100
$\Delta(\Delta_r G^\circ) = -RT\ln 100 = \qty{-11.4}{kJ/mol}$.
\end{solution}

\begin{solution}{exo:b2:equilibrium-shifts:2}
$\ln K^\circ(400) = \ln(5.4 \times 10^5) + (91\,880/8.314)(1/400 -
1/298.15) = 13.20 - 9.44 = 3.76$، $K^\circ \approx 43$. وتعطي الجداول 35.6:
فعلى مدى \qty{100}{K} يخطئ تقريب ثبات $\Delta_r H^\circ$ منذ الآن
بنحو 20\,\% (إذ يصير التفاعل أشد نشرًا للحرارة مع ارتفاع $T$).
\end{solution}

\begin{solution}{exo:b2:equilibrium-shifts:3}
(أ) $n = 1$، $r = 0$، $\varphi = 2$: $v = 1$ (فضغط البخار دالة
في $T$). (ب) $v = 3 - 1 + 2 - 3 = 1$. (ج) $v = 3 - 1 + 2 - 1 = 3$. (د) أربعة
أنواع، وتفاعل واحد، وطوران (الكربون والغاز): $v = 4 - 1 + 2 - 2 =
3$؛ وإذا لم يُصنع الغاز إلا من الكربون وبخار الماء، كان $n(\ce{CO}) = n(\ce{H2})$
في الغاز، و$s = 1$ و$v = 2$.
\end{solution}

\begin{solution}{exo:b2:equilibrium-shifts:4}
التسخين: الاتجاه \omterm{def:b2:reaction-enthalpy:exothermic}{الماص للحرارة}، نحو \ce{NO2} (المباشر). الضغط:
نحو جزيئات غازية أقل، أي \ce{N2O4} (المعاكس). والأرغون عند حجم ثابت:
لا تغيير، فالضغوط الجزئية للغازات المتفاعلة لا تتغير.
\end{solution}

\begin{solution}{exo:b2:equilibrium-shifts:5}
عند \qty{298}{K}: $\Delta_r G^\circ = 180.6 - 298.15 \times 0.0248 =
\qty{173.2}{kJ/mol}$، $K^\circ = \num{4.5e-31}$. وعند \qty{2000}{K}: $\Delta_r
G^\circ = 180.6 - 2000 \times 0.0248 = \qty{131.0}{kJ/mol}$، $K^\circ =
\num{3.8e-4}$. ومع $\Delta\nu_{\text{غاز}} = 0$، $x(\ce{NO})^2 = K^\circ x(\ce{N2})
x(\ce{O2})$: $x(\ce{NO}) = \sqrt{3.8 \times 10^{-4} \times 0.78 \times 0.21} =
\num{7.9e-3}$، أي واحد في المئة تقريبًا. والمحركات تحترق عند درجات حرارة
كهذه؛ ومع تبرد الغازات ينخفض $K^\circ$ انخفاضًا هائلًا، لكن تفكك \ce{NO}
بطيء جدًا: فيتجمد التوازن الساخن في غازات العادم (إلى أن يزيله
حفاز).
\end{solution}

\begin{solution}{exo:b2:equilibrium-shifts:6}
مع $\xi$ مول من الإستر وكمية مادة كلية ثابتة، $Q =
\xi^2/[(1 - \xi)(n_B - \xi)]$. من أجل $n_B = 1.0$: $\xi^2/(1 - \xi)^2 = 4$، $\xi
= \qty{0.67}{mol}$. ومن أجل $n_B = 3.0$: $3\xi^2 - 16\xi + 12 = 0$، $\xi =
\qty{0.90}{mol}$. ونزع الماء يُبقي $Q$ أدنى من $K^\circ$ أيًّا كانت $\xi$: فيستمر
التفاعل في الاتجاه المباشر حتى يُستهلك الحمض.
\end{solution}

\begin{solution}{exo:b2:equilibrium-shifts:7}
دون أرغون: كميات المادة $1 - \alpha$ و$2\alpha$، والمجموع $1 + \alpha$؛
$Q = 4\alpha^2/(1 - \alpha^2) = 0.148$، $\alpha = \sqrt{0.148/4.148} = 0.189$.
ومع \qty{1}{mol} من الأرغون عند \qty{1}{bar}: المجموع $2 + \alpha$، و$Q =
4\alpha^2/[(1 - \alpha)(2 + \alpha)] = 0.148$، ومنه $4.148\alpha^2 + 0.148\alpha
- 0.296 = 0$، $\alpha = 0.250$: فالتخفيف عند ضغط كلي ثابت يخفض
الضغوط الجزئية، وهو ما يحابي الجانب ذا الجزيئات الغازية الأكثر. وعند
حجم ثابت لا تتغير الضغوط الجزئية للغازات المتفاعلة:
فيبقى $\alpha$ مساويًا 0.189.
\end{solution}

\begin{solution}{exo:b2:equilibrium-shifts:8}
الجسم الصلب وحده: $n = 3$، $r = 1$، $s = 1$ ($p(\ce{NH3}) = p(\ce{HCl})$)،
$\varphi = 2$: $v = 3 - 1 - 1 + 2 - 2 = 1$؛ اختر $T$ فيتحدد
الضغط («ضغط التفكك»). ومع إضافة الأمونيا، $s = 0$ و$v = 2$:
يمكن اختيار $T$ وضغط واحد؛ والجداء $p(\ce{NH3})\,p(\ce{HCl})$ تحدده
$T$.
\end{solution}

\begin{solution}{exo:b2:equilibrium-shifts:9}
$\Delta_r H^\circ = -R\ln(K_2/K_1)/(1/T_2 - 1/T_1) = -8.314 \times
\ln(0.0173)/(1/600 - 1/500) = \qty{-101}{kJ/mol}$، أي أكثر سالبية منه عند
\qty{298}{K} (\qty{-91.9}{kJ/mol})، كما تنبأ \omterm{thm:b2:reaction-enthalpy:kirchhoff}{قانون كيرشوف}
($\Delta_r C_p^\circ < 0$).
\end{solution}

\begin{solution}{exo:b2:equilibrium-shifts:10}
$Q = \dfrac{n_{\ce{NH3}}^2\,n_{\text{كلي}}^2}{n_{\ce{N2}}n_{\ce{H2}}^3}
\left(\dfrac{p^\circ}{p}\right)^2$. وعند $T$ و$p$ ثابتين:
$\partial\ln Q/\partial n_{\ce{N2}} = -1/n_{\ce{N2}} + 2/n_{\text{كلي}}$، وهو موجب
حين $n_{\ce{N2}} > n_{\text{كلي}}/2$، أي $x(\ce{N2}) > 1/2$. فيتجاوز $Q$
القيمة $K^\circ$ وينزاح التوازن في الاتجاه المعاكس. وعند حجم ثابت،
$\partial\ln Q/\partial n_{\ce{N2}} = -1/n_{\ce{N2}} < 0$: دائمًا في الاتجاه المباشر.
\end{solution}

\begin{solution}{exo:b2:equilibrium-shifts:11}
$K^\circ = \exp(-1620/(8.314 \times 1100)) = 0.84$: $p(\ce{CO2}) =
\qty{0.84}{bar}$ عند التوازن، أي $n = pV/RT = 0.84 \times 10^5 \times
0.0100/(8.314 \times 1100) = \qty{0.092}{mol}$ من الغاز. ومع \qty{0.050}{mol}
من الحجر الجيري يتفكك كله ($p = \qty{0.46}{bar} < K^\circ p^\circ$):
فلا توازن، بل قطع للتوازن. ومع \qty{0.200}{mol} يحدث التوازن:
\qty{0.092}{mol} من \ce{CO2} ومن \ce{CaO}، ويبقى \qty{0.108}{mol} من \ce{CaCO3}.
\end{solution}

\begin{solution}{exo:b2:equilibrium-shifts:12}
تخرج الطبقة 1 عند نحو \qty{890}{K} وتحويل قدره 68\,\%، والطبقة 2 عند نحو
\qty{785}{K} و92\,\%، والطبقة 3 عند نحو \qty{725}{K} و98\,\%. ففي طبقة
كظومة واحدة يسخن الغاز وهو يتفاعل حتى يلاقي منحنى
التوازن، الذي لا يسمح عند تلك الدرجة إلا بنحو 68\,\%. ولا يمكن إبقاء الحفاز
باردًا طوال الوقت لأنه لا يعمل إلا ساخنًا؛ والتبريد
بين الطبقات يحفظ السرعة ويرفع حد التوازن. والطبقة الأخيرة
يحدّها التوازن عند أدنى درجة حرارة يعمل فيها الحفاز
بعدُ (ولهذا تمتص المصانع ثالث الأكسيد ثم تمرر الغاز مرة أخرى على
حفاز).
\end{solution}

\begin{solution}{pb:b2:equilibrium-shifts:1}
\textbf{1.} $\Delta_r H^\circ = \qty{-91.88}{kJ/mol}$، $\Delta_r S^\circ =
2(192.77) - 191.61 - 3(130.68) = \qty{-198.1}{J/(K.mol)}$، $\Delta_r G^\circ =
\qty{-32.8}{kJ/mol}$، $K^\circ = \num{5.6e5}$.
\textbf{2.} $K^\circ(700) = \exp(-54\,380/(8.314 \times 700)) = \num{8.8e-5}$؛
$K^\circ(800) = \exp(-77\,320/(8.314 \times 800)) = \num{8.9e-6}$.
\textbf{3.} $\Delta_r H^\circ = -R\ln(K_{800}/K_{700})/(1/800 - 1/700)$،
أي \qty{-106}{kJ/mol}، وهي أكثر سالبية منها عند \qty{298}{K}.
\textbf{4.} $\ln K^\circ(723) = \ln(8.75 \times 10^{-5}) + 12\,777 \times
(1/723.15 - 1/700)$، مع $12\,777 = 106\,230/8.314$: $K^\circ =
\num{4.9e-5}$.
\textbf{5.} $\Delta_r G^\circ(723) = -91.88 + 723.15 \times 0.1981 =
\qty{51.4}{kJ/mol}$، $K^\circ = \num{1.9e-4}$: أكبر بأربع مرات مما ينبغي، لأن
$\Delta_r H^\circ$ و$\Delta_r S^\circ$ يتغيران على مدى \qty{425}{K}
($\Delta_r C_p^\circ \approx \qty{-44}{J/(K.mol)}$).
\textbf{6.} \ce{N2} $1 - \xi$، \ce{H2} $3 - 3\xi$، \ce{NH3} $2\xi$، والمجموع $4 -
2\xi$.
\textbf{7.} $x = 2\xi/(4 - 2\xi)$، ومنه $1 - x = (4 - 4\xi)/(4 - 2\xi)$؛ فكسر
النيتروجين $(1 - \xi)/(4 - 2\xi) = (1 - x)/4$ وكسر الهيدروجين
أكبر منه بثلاث مرات.
\textbf{8.} لدينا
\[
  K^\circ = \frac{x^2}{\frac{1-x}{4}\left(\frac{3(1-x)}{4}\right)^3}
  \left(\frac{p^\circ}{p}\right)^2 = \frac{256}{27}\,\frac{x^2}{(1 - x)^4}
  \left(\frac{p^\circ}{p}\right)^2 ;
\]
ثم نأخذ الجذر التربيعي.
\textbf{9.} $a = (\sqrt{27}/16) \times \sqrt{4.9 \times 10^{-5}} \times 1 =
\num{2.27e-3}$؛ و$x/(1 - x)^2 = a$ تعطي $x = 0.0023$.
\textbf{10.} $a = 0.455$؛ $ax^2 - (2a + 1)x + a = 0$: $x = 0.25$.
\textbf{11.} $n = 3$، $r = 1$، $s = 1$ (النسبة 1 : 3 يحفظها التفاعل)،
$\varphi = 1$: $v = 2$. فمتى اختيرت $T$ و$p$ لم يبق شيء حر.
\textbf{12.} الضغط: في الاتجاه المباشر ($\Delta\nu_{\text{غاز}} = -2$). ودرجة الحرارة:
في الاتجاه المعاكس (\omterm{def:b2:reaction-enthalpy:exothermic}{ناشر للحرارة}).
\textbf{13.} الغاز الخامل عند ضغط كلي ثابت يخفض الضغوط الجزئية
للغازات المتفاعلة، كما يفعل انخفاض الضغط: في الاتجاه المعاكس.
\textbf{14.} $\partial\ln Q/\partial n_{\ce{N2}} = (-1/0.30 + 2)/n_{\text{كلي}} < 0$:
فينخفض $Q$ تحت $K^\circ$، في الاتجاه المباشر.
\textbf{15.} لا: عند حجم ثابت $\partial\ln Q/\partial n_{\ce{N2}} =
-1/n_{\ce{N2}} < 0$ أيًّا كان التركيب: في الاتجاه المباشر.
\textbf{16.} وتكاد الأمونيا تنعدم، فيكون $Q$ أدنى بكثير من $K^\circ$ عند المدخل:
فيدخل الغاز بعيدًا عن التوازن ويتفاعل في الاتجاه المباشر.
\textbf{17.} يمكث الغاز على الحفاز زمنًا أقصر من أن يبلغ
التوازن؛ وتماسٌّ أطول يكلّف محوِّلًا أكبر وأغلى
لكسب صغير.
\textbf{18.} $x = 2\xi/(4 - 2\xi) = 0.18$ تعطي $\xi = \qty{0.305}{mol}$ لكل
مول من \ce{N2} المغذّى: \omterm{def:b2:equilibrium-shifts:single-pass}{التحويل في مرور واحد} للهيدروجين (وللنيتروجين)
$0.305$، أي 31\,\%.
\textbf{19.} في التشغيل المستقر تتحول التغذية الطازجة كلها في النهاية:
فيخرج \qty{2.00}{mol} من الأمونيا في وحدة الزمن؛ ويجب أن يمر في المحوِّل $1.00/
0.305 = \qty{3.3}{mol}$ من \ce{N2} في وحدة الزمن، ويُعاد تدوير الباقي.
\textbf{20.} الأرغون (والميثان) الداخلان مع التغذية لا يتفاعلان أبدًا ولا
يُزالان مع الأمونيا السائلة: فبلا \omterm{def:b2:equilibrium-shifts:single-pass}{تنفيس} يتراكمان في الحلقة
ويخفضان الضغوط الجزئية للمتفاعلات.
\textbf{21.} عند \qty{500}{K} يكون حفاز الحديد بطيئًا جدًا: فالتوازن
يكون مناسبًا لكنه لا يُقارَب أبدًا في زمن صناعي.
\textbf{22.} عند \qty{723}{K} و\qty{200}{bar}، وبتغذية ستوكيومترية وغازات
مثالية: $\boldsymbol{x(\ce{NH3}) \approx 0.25}$.
\end{solution}
